\section{Case Study 3: Qwen 3}

最后一个案例是 Qwen 3。课堂给它的定位是“更晚出现、但更系统地整合了前面这些经验”的一个开放模型尝试。

\begin{figure}[H]
\centering
\includegraphics[width=0.95\textwidth]{slides-images/slide-061.jpg}
\caption{Qwen 3：更晚期的开源 reasoning model 尝试}
\end{figure}

\subsection{SFT + reasoning RL + 组合策略}

Qwen 3 的 recipe 里有几个很清楚的点：
\begin{itemize}
    \item 对训练数据做难度过滤，优先保留模型最开始做不好的题。
    \item 清理掉模型本来就能轻松做对的题，避免把算力浪费在简单样本上。
    \item 对 CoT 的质量做人工或半人工过滤，区分“猜对”和“真推对”。
    \item 用很小的一批样本也能启动 reasoning RL。
\end{itemize}

之后再进入更成熟的 post-training 流程：先 reasoning RL，再做常规 RLHF，最后配合蒸馏和 test-time scaling。

\begin{figure}[H]
\centering
\includegraphics[width=0.95\textwidth]{slides-images/slide-063.jpg}
\caption{Qwen 3 的 SFT + reasoning RL：数据过滤和 GRPO 是关键部分}
\end{figure}

\begin{figure}[H]
\centering
\includegraphics[width=0.95\textwidth]{slides-images/slide-064.jpg}
\caption{thinking mode fusion：用标签和 early stop 控制 CoT 长度}
\end{figure}

\begin{figure}[H]
\centering
\includegraphics[width=0.95\textwidth]{slides-images/slide-065.jpg}
\caption{test-time scaling：除了训练，推理时的计算也可以继续放大效果}
\end{figure}

\begin{figure}[H]
\centering
\includegraphics[width=0.95\textwidth]{slides-images/slide-066.jpg}
\caption{整体阶段组合：math / STEM 往往在更一般的 RLHF 阶段里会略有回落}
\end{figure}

